\pdfoutput=1
\documentclass[11pt]{article}
\usepackage[a4paper,margin=2.6cm]{geometry}
\usepackage{amsmath,amssymb,amsthm,mathtools}
\usepackage{booktabs}
\usepackage[breakable,skins]{tcolorbox}
\usepackage[colorlinks=true,linkcolor=blue!50!black,citecolor=blue!50!black,urlcolor=blue!50!black]{hyperref}
\usepackage{enumitem}

\newtheorem{theorem}{Theorem}[section]
\newtheorem{proposition}[theorem]{Proposition}
\newtheorem{lemma}[theorem]{Lemma}
\newtheorem{corollary}[theorem]{Corollary}
\theoremstyle{definition}
\newtheorem{definition}[theorem]{Definition}
\newtheorem{inputthm}[theorem]{External input}
\newtheorem{remark}[theorem]{Remark}
\newtheorem{exampleblock}[theorem]{Example}

\newtheorem*{theoremA}{Theorem A}
\newtheorem*{corollaryB}{Corollary B}
\newtheorem*{corollaryC}{Corollary C}

\tcolorboxenvironment{proof}{breakable,enhanced,colback=gray!3,colframe=gray!45,boxrule=0.4pt,arc=1pt,left=6pt,right=6pt,top=4pt,bottom=4pt}
\tcolorboxenvironment{exampleblock}{breakable,enhanced,colback=blue!2,colframe=blue!30,boxrule=0.4pt,arc=1pt,left=6pt,right=6pt,top=4pt,bottom=4pt}

\newcommand{\g}{\mathfrak g}
\newcommand{\h}{\mathfrak t}
\newcommand{\Cs}{\mathcal C}
\newcommand{\Oo}{\mathcal O}
\newcommand{\rk}{\operatorname{rk}}
\newcommand{\Sf}{\mathbb S}
\newcommand{\hv}{h^\vee}
\newcommand{\Sg}{\mathfrak S}
\newcommand{\Asxx}{A^{\mathrm{SXX}}}
\DeclareMathOperator{\Sym}{Sym}
\DeclareMathOperator{\Hom}{Hom}
\DeclareMathOperator{\GL}{GL}
\DeclareMathOperator{\SL}{SL}
\DeclareMathOperator{\Sp}{Sp}
\DeclareMathOperator{\Pf}{Pf}
\DeclareMathOperator{\sgn}{sgn}

\title{Gram determinants on commuting schemes\\ and the Chevalley restriction map}
\author{Botao Li}
\date{October 2026}

\begin{document}
\maketitle

\begin{abstract}
Let $\g$ be a simple Lie algebra over an algebraically closed field of characteristic zero, $\h\subset\g$ a Cartan subalgebra with Weyl group $W$, and $B$ an invariant form. Chen and Ng\^o conjectured that restriction to $\h^d$ identifies the invariant functions on the scheme of commuting $d$-tuples in $\g$ with $W$-invariant functions on $\h^d$. We determine the $\SL_d$-invariant part of degree $2d$ of this invariant ring: it is one-dimensional when $d$ is smaller than the dual Coxeter number $\hv$, spanned by the Gram determinant $\det\bigl(B(x_i,x_j)\bigr)$, and zero for $d\ge\hv$. The proof transfers the Cachazo--Douglas--Seiberg--Witten theorem on the classical chiral ring of $\mathcal N=1$ super-Yang--Mills theory, in which the gaugino bilinear $S$ satisfies $S^{\hv-1}\neq0$ and $S^{\hv}=0$, to commuting tuples. For $\rk\g<d<\hv$ the Gram determinant is a nonzero nilpotent invariant with zero restriction to $\h^d$, so the restriction map has a nonzero kernel. This range is nonempty for $\mathfrak{so}_n$ with $n\ge7$ (smallest case $\mathfrak{so}_7$, $d=4$) and for every exceptional type (smallest case $\g_2$, $d=3$), and empty in types A and C, where the conjecture is a theorem. The case $\mathfrak{so}_7$, $d=4$ contradicts a published theorem of Song, Xia and Xu; we show that their Lemma~5.2 fails for every $d\ge2$ and locate the step of its proof that breaks. In the super-commuting setting of Chang, the same component is an irreducible $\mathfrak{gl}(3|2)$-module of the classical $Q$-cohomology of $\mathcal N=4$ super-Yang--Mills theory whose purely odd vector is $S^d$; for $\rk\g<d<\hv$ these gaugino powers are the non-Coulomb monotone classes of Kibe, van Leuven and Mathieson, and the Gram determinant is their counterpart on the purely even commuting scheme. With the count of Chang and Lin, the $\mathfrak{so}_7$ component at $d=4$ contains the non-Cartan class in the first classical mismatch between $\mathfrak{so}_7$ and $\mathfrak{sp}_6$.
\end{abstract}

\section{Introduction}

Let $k$ be an algebraically closed field of characteristic zero, $G$ a reductive group over $k$ with Lie algebra $\g$, $\h\subset\g$ a Cartan subalgebra and $W=N_G(\h)/Z_G(\h)$. For $d\ge1$ the commuting scheme $\Cs^d_\g\subset\g^d$ is the scheme-theoretic zero fibre of the commutator map
\[
\g^d\longrightarrow\prod_{i<j}\g,\qquad (x_1,\dots,x_d)\longmapsto\bigl([x_i,x_j]\bigr)_{i<j},
\]
so $k[\Cs^d_\g]=k[\g^d]/I_d$, where $I_d$ is generated by the coordinate functions of all $[x_i,x_j]$. Restriction along $\h^d\subset\Cs^d_\g$ defines
\[
\Phi:\ k[\Cs^d_\g]^G\longrightarrow k[\h^d]^W .
\]
For $d=1$ this is the Chevalley restriction isomorphism. Chen and Ng\^o \cite[Conj.~3.1]{CN20} conjectured that $\Phi$ is an isomorphism for every reductive $G$ and every $d$, equivalently that $\Cs^d_\g/\!\!/G$ is reduced and normal; their weaker Conjecture~3.2 (existence of a spectral data morphism) already forces $\Phi$ to be injective \cite[\S3]{CN20}. The conjecture is a theorem for $\GL_n$ (Vaccarino \cite{Va07}, Domokos \cite{Do09}; Chen and Ng\^o prove Conjecture~3.2 for $\GL_n$ in \cite[Thm.~3.3]{CN20}) and for $\Sp_{2n}$ (Chen--Ng\^o \cite[Thm.~4.1]{CN21}); Song, Xia and Xu \cite[Thm.~1.2]{SXX23} state it for $O_n$ and $SO_n$. For $d=2$ it follows for every connected reductive group from Li--Nadler--Yun \cite[Thm.~1.4.11]{LNY24}, stated there under their Ansatz~1.4.2, and from the normality of $\Cs^2_\g$ proved by Sempliner and Yang \cite{SY26} together with \cite[\S3]{CN20}; for $\g_2$ see also Goyal \cite{Go26}. For reduced schemes, $\h^d/\!\!/W\to(\Cs^d_{\g})_{\mathrm{red}}/\!\!/G$ is an isomorphism for types A, B, C, D and $G_2$ by Hunziker \cite{Hu97}, and for every reductive $G$ it is finite and bijective on $k$-points \cite[Thms.~6.2, 6.3]{Hu97}, building on Richardson \cite{Ri88}.

We study the Gram determinant. Fix a nondegenerate invariant symmetric form $B$ on $\g$ and set
\[
T(x)=\bigl(B(x_i,x_j)\bigr)_{1\le i,j\le d},\qquad \det T\in k[\g^d]^{G}_{2d}.
\]
Under $\GL_d$ acting on $(x_1,\dots,x_d)$, $\det T$ transforms by $\det^2$, so it lies in the $\SL_d$-invariant part of degree $2d$. When $d>\rk\g$, any $d$ vectors of $\h$ are linearly dependent and $\det T|_{\h^d}=0$. The question is whether $\det T$ lies in $I_d$.

The answer comes from the classical chiral ring of four-dimensional $\mathcal N=1$ super-Yang--Mills theory. Cachazo, Douglas, Seiberg and Witten \cite{CDSW} considered the supercommutative algebra $A$ generated by two odd copies $X,Y$ of $\g$ subject to $\{X,X\}=\{X,Y\}=\{Y,Y\}=0$, and the invariant $S=B(X,Y)$. Their conjecture, proved for every simple $\g$ \cite{CDSW,Wi03,EK03,Et03,Ku08,CF23}, states that $A^\g$ is generated by $S$, that $S^{\hv}=0$ and that $S^{\hv-1}\neq0$, where $\hv$ is the dual Coxeter number. The substitution $x_i=Xu_i+Yv_i$, with odd auxiliary variables $u_i,v_i$, produces commuting $d$-tuples on which $\det T$ takes the value $(-1)^d(d+1)!\,S^d\prod_iu_iv_i$. Super Schur--Weyl duality identifies the $\SL_d$-invariant part of $k[\Cs^d_\g]^G$ of degree $2d$ with the $\SL_2$-invariant part of $A^\g$ of degree $2d$, which is spanned by $S^d$.

\begin{theoremA}[Theorem~\ref{thm:A}]
Let $\g$ be simple and $d\ge1$. The $\SL_d$-invariant subspace of $k[\Cs^d_\g]^G$ of degree $2d$ has dimension $1$ if $d<\hv$ and $0$ if $d\ge\hv$. For $d<\hv$ it is spanned by the class of $\det T$.
\end{theoremA}

\begin{corollaryB}[Corollary~\ref{cor:B}]
If $\rk\g<d<\hv$, then the class of $\det T$ is a nonzero nilpotent element of $k[\Cs^d_\g]^G$ with $\Phi(\det T)=0$. Hence $\ker\Phi\neq0$, the nilradical of $k[\Cs^d_\g]^G$ is nonzero, and Conjectures~3.1 and~3.2 of \cite{CN20} fail for $(\g,d)$.
\end{corollaryB}

The range $\rk\g<d<\hv$ is empty for types A and C, in agreement with the theorems of Vaccarino, Domokos and Chen--Ng\^o, and nonempty for $B_n$ ($n\ge3$), $D_n$ ($n\ge4$) and every exceptional type (Table~\ref{tab:windows}). The smallest case is $\g_2$ with $d=3$; there $\det T$ spans $\ker\Phi$ in all degrees at most $10$, and $(\det T)^2\in I_3$ (Example~\ref{ex:g2}).

\begin{corollaryC}[Corollary~\ref{cor:SXX}]
For $n\ge7$, \cite[Thm.~1.2(1), Cor.~1.3]{SXX23} fail for $G=O_n$ and $G=SO_n$; the smallest case is $n=7$, $d=4$. \cite[Lemma~5.2]{SXX23} fails for every $d\ge2$.
\end{corollaryC}

Section~\ref{sec:sxx} proves the statement about \cite[Lemma~5.2]{SXX23} directly: on the hypersurface where that lemma looks for solutions, the $2d\times2d$ matrix it uses has rank at most $3$. Section~\ref{sec:N4} treats the super version of Theorem~A, which concerns the classical $Q$-cohomology of $\mathcal N=4$ super-Yang--Mills theory.\footnote{By \cite{CY13} and \cite[Prop.~4]{Ch26}, the classical $Q$-cohomology of top degree is the invariant ring of a super-commuting scheme with three even and two odd directions, and its purely even part is $k[\Cs^3_\g]^G$. For $\g_2$ the class of $\det T$ lies in Chang's multidegree $n=(0,0,2,2,2)$ (charges $(J_1,J_2,q_1,q_2,q_3)=(0,0,2,2,2)$ in the conventions of \cite[\S5.2]{Ch26}) and has zero Cartan restriction.}

Kibe, van Leuven and Mathieson \cite[eq.~(1)]{KvLM26} identified the powers $S^m$, $\rk\g<m<\hv$, of the gaugino bilinear as non-Coulomb monotone classes of this $Q$-cohomology, nonzero by External input~\ref{in:CDSW}; they check the nonvanishing directly for $\mathfrak{so}_7,\dots,\mathfrak{so}_{12}$ \cite[\S4.5]{KvLM26}, relate $S^4$ for $\mathfrak{so}_7$ to the class of \cite{CL25,CE25}, and construct a fortuitous partner of $S^4$ under the one-loop supercharge \cite[\S5]{KvLM26}. Their classes are the purely odd vectors of the modules of Corollary~\ref{cor:N4} (Remark~\ref{rem:N4}(d)). The present note adds the passage to the purely even commuting scheme: by Theorem~A the gaugino power $S^d$ on $\Cs_{\Pi k^2}$ and the Gram determinant on $\Cs^d_\g$ have the same multiplicity space $L_d$. This gives Corollaries~B and~C, the full $\Sf_{(2^d)}$-component for every $E$ (Corollary~\ref{cor:N4}), and injectivity of the super Cartan restriction for $\mathfrak{sp}_{2N}$ (Corollary~\ref{cor:spE}).

\section{Setting and statements}\label{sec:setting}

\begin{definition}[Commuting superschemes]\label{def:CE}
Let $E$ be a finite-dimensional super vector space over $k$. The space $\Hom(E,\g)$ has coordinate superalgebra $\Sym(E\otimes\g^*)$, the supersymmetric algebra. For $e,e'\in E$ and $\xi\in\g^*$ let $\rho_E(e\wedge e'\otimes\xi)\in\Sym^2(E\otimes\g^*)$ be the function $\theta\mapsto\xi([\theta(e),\theta(e')])$; it depends on $e\wedge e'$ in the super exterior square $\Lambda^2E$. Set
\[
\Oo(\Cs_E)=\Sym(E\otimes\g^*)\big/J_E,\qquad J_E=\bigl(\rho_E(\Lambda^2E\otimes\g^*)\bigr).
\]
Its points over a supercommutative $k$-algebra $R$ are the even elements $\theta$ of $\Hom(E,\g)\otimes R$ with $[\theta(e),\theta(e')]=0$ for all $e,e'\in E$. The group $\GL(E)\times G$ acts. For $E=k^d$ this is the commuting scheme, $\Oo(\Cs_{k^d})=k[\Cs^d_\g]$ \cite[\S3]{CN20}; for $E=k^{d_B|d_F}$ it is the super-commuting scheme $\Cs^{d_B|d_F}_\g$ of Chang~\cite[\S2]{Ch26}. The Cartan restriction is
\[
\Phi_E:\Oo(\Cs_E)^G\longrightarrow\Oo(\Hom(E,\h))^W .
\]
\end{definition}

\begin{inputthm}[{CDSW conjecture. Part (i) for every simple $\g$: Kumar \cite{Ku08}. Parts (i)--(iii) for classical $\g$: \cite{CDSW,Wi03} (for (iii) in types B and D, \cite{Wi03} uses instantons), algebraically \cite{Et03}; for $\g_2$: \cite{EK03}; for $F_4,E_6,E_7,E_8$: \cite[Prop.~2, Thm.~1, Cor.~2]{CF23}}]\label{in:CDSW}
Let $\g$ be simple with dual Coxeter number $\hv$, $(e_a)$ a basis of $\g$ and $X^a,Y^a$ odd generators. Put $X=\sum_ae_a\otimes X^a$, $Y=\sum_ae_a\otimes Y^a$, and let $A$ be the quotient of $\Lambda(X^a,Y^a)$ by the coordinates of
\[
\{X,X\}=\{X,Y\}=\{Y,Y\}=0,\qquad\text{where}\quad \{X,Y\}:=\textstyle\sum_{a,b}[e_a,e_b]\otimes X^aY^b
\]
and similarly for $\{X,X\}$, $\{Y,Y\}$; equivalently, $A=\Lambda(\g\oplus\g)/I$ with $I$ generated by the copies of $\g$ in bidegrees $(2,0)$, $(1,1)$, $(0,2)$ \cite[\S1]{EK03}. Let $S=\sum_{a,b}B(e_a,e_b)X^aY^b$. Then \cite[Conj.~1.1]{EK03}
\[
\text{(i)}\ \ A^\g\ \text{is generated by } S,\qquad \text{(ii)}\ \ S^{\hv}=0,\qquad \text{(iii)}\ \ S^{\hv-1}\neq0 .
\]
The algebra $A$, the element $S$ and the action of $G$ are defined over $\mathbb Q$, so (i)--(iii) hold over every field of characteristic zero. In the notation of Definition~\ref{def:CE}, $A=\Oo(\Cs_{\Pi k^2})$, where $\Pi k^2$ is the purely odd space of dimension $0|2$ and $\g^*$ is identified with $\g$ through $B$.
\end{inputthm}

\begin{theorem}[Gram component]\label{thm:A}
Let $\g$ be simple, $G$ connected and $d\ge1$, and let $\Sf_\lambda$ be the super Schur functor (External input~\ref{in:schur}). There is a vector space $L_d$, independent of $E$, such that for every finite-dimensional super vector space $E$ the $\Sf_{(2^d)}(E)$-isotypic component of the $\GL(E)$-module $\Oo(\Cs_E)^G_{2d}$ is
\[
\Oo(\Cs_E)^G_{(2^d)}\ \cong\ \Sf_{(2^d)}(E)\otimes L_d,\qquad
\dim L_d=\begin{cases}1,& d<\hv,\\ 0,& d\ge\hv,\end{cases}
\]
where $(2^d)=(2,\dots,2)$ has $d$ parts. For $E=k^d$ this component is the $\SL_d$-invariant subspace of degree $2d$ of $k[\Cs^d_\g]^G$, and for $d<\hv$ it is spanned by the class of $\det T$.
\end{theorem}

\begin{corollary}\label{cor:B}
Let $\g$ be simple with $\rk\g<d<\hv$. Then the class of $\det T$ in $k[\Cs^d_\g]^G$ is nonzero and nilpotent, and $\Phi(\det T)=0$. Consequently $\ker\Phi\neq0$, the nilradical of $k[\Cs^d_\g]^G$ is nonzero, and Conjectures~3.1 and~3.2 of \cite{CN20} fail for $(\g,d)$.
\end{corollary}

\begin{table}[ht]
\centering\small
\begin{tabular}{lccl}
\toprule
type & $\rk$ & $\hv$ & $d$ with $\rk<d<\hv$\\
\midrule
$A_n$ & $n$ & $n+1$ & --\\
$B_n$, $n\ge2$ & $n$ & $2n-1$ & $n+1\le d\le 2n-2$\quad ($n\ge3$)\\
$C_n$, $n\ge3$ & $n$ & $n+1$ & --\\
$D_n$, $n\ge4$ & $n$ & $2n-2$ & $n+1\le d\le 2n-3$\\
$G_2$ & $2$ & $4$ & $d=3$\\
$F_4$ & $4$ & $9$ & $5\le d\le8$\\
$E_6$ & $6$ & $12$ & $7\le d\le11$\\
$E_7$ & $7$ & $18$ & $8\le d\le17$\\
$E_8$ & $8$ & $30$ & $9\le d\le29$\\
\bottomrule
\end{tabular}
\caption{Pairs $(\g,d)$ covered by Corollary~\ref{cor:B}. For $\mathfrak{so}_n$: $n=2m+1\ge7$ with $m+1\le d\le2m-2$, and $n=2m\ge8$ with $m+1\le d\le2m-3$.}\label{tab:windows}
\end{table}

\section{The transfer}\label{sec:transfer}

\begin{lemma}[Grassmann Gram determinant]\label{lem:grass}
Let $\Lambda=\Lambda[u_1,\dots,u_d,v_1,\dots,v_d]$ be the exterior algebra on odd generators and $N_{ij}=u_iv_j+u_jv_i\in\Lambda_{\bar0}$. Then
\[
\det(N_{ij})_{1\le i,j\le d}=(d+1)!\;u_1v_1u_2v_2\cdots u_dv_d .
\]
\end{lemma}
\begin{proof}
Adjoin odd variables $\theta_i,\bar\theta_i$ and let $\int$ extract the coefficient of $\Theta=\bar\theta_1\theta_1\cdots\bar\theta_d\theta_d$. For a matrix $M$ with entries in the commutative ring $\Lambda_{\bar0}$, the top component of $\exp(\sum_{i,j}\bar\theta_iM_{ij}\theta_j)$ is
\[
\tfrac1{d!}\Bigl(\sum_{i,j}\bar\theta_iM_{ij}\theta_j\Bigr)^d=\sum_{\sigma\in\Sg_d}\prod_{r}M_{r\sigma(r)}\,\bar\theta_r\theta_{\sigma(r)}=\det M\;\Theta ,
\]
since reordering the odd $\theta$'s among the even pairs $\bar\theta_r\theta_{\sigma(r)}$ produces $\sgn\sigma$. Hence $\det M=\int\exp(\bar\theta M\theta)$. For $M=N$, with the even elements
\[
\alpha=\sum_i\bar\theta_iu_i,\quad \beta=\sum_jv_j\theta_j,\quad \gamma=\sum_i\bar\theta_iv_i,\quad \varepsilon=\sum_ju_j\theta_j,
\]
one has $\bar\theta_iu_jv_i\theta_j=-\bar\theta_iv_i\,u_j\theta_j$ and therefore $\sum_{i,j}\bar\theta_iN_{ij}\theta_j=\alpha\beta-\gamma\varepsilon$. As $\alpha,\beta,\gamma,\varepsilon$ commute,
\[
\det N=\int\sum_{r+s=d}\frac{(-1)^s}{r!\,s!}\,\alpha^r\gamma^s\,\beta^r\varepsilon^s .
\]
Expanding the powers,
\[
\alpha^r\gamma^s=r!\,s!\sum_{I\sqcup J=[d],\,|I|=r}\ \prod_{i\in I}\bar\theta_iu_i\prod_{j\in J}\bar\theta_jv_j,\qquad
\beta^r\varepsilon^s=r!\,s!\sum_{K\sqcup L=[d],\,|K|=r}\ \prod_{i\in K}v_i\theta_i\prod_{j\in L}u_j\theta_j .
\]
A product is nonzero exactly when every $u_m$ and every $v_m$ occurs once, i.e.\ when $L=J$ and $K=I$. Then
\[
\bar\theta_iu_i\,v_i\theta_i=\bar\theta_i\theta_i\,u_iv_i,\qquad \bar\theta_jv_j\,u_j\theta_j=-\bar\theta_j\theta_j\,u_jv_j,
\]
so each surviving term equals $(-1)^s\,\Theta\,\prod_mu_mv_m$. Collecting the coefficients,
\[
\det N=\sum_{r=0}^d\frac{(-1)^s}{r!\,s!}\,(r!\,s!)^2\binom dr(-1)^s\prod_mu_mv_m=\sum_{r=0}^dd!\,\prod_mu_mv_m=(d+1)!\prod_mu_mv_m .
\]
\end{proof}

\begin{proposition}[Transfer]\label{prop:transfer}
Let $\g$ be simple, $A$ as in External input~\ref{in:CDSW}, $\Lambda$ as in Lemma~\ref{lem:grass}, and $R_d=(A\otimes\Lambda)_{\bar0}$, a commutative $k$-algebra. Put
\[
x_i=Xu_i+Yv_i=\sum_ae_a\otimes\bigl(X^au_i+Y^av_i\bigr)\in\g\otimes R_d,\qquad 1\le i\le d.
\]
Then:
\begin{enumerate}[label=(\roman*)]
\item $[x_i,x_j]=0$ for all $i,j$, so $(x_1,\dots,x_d)\in\Cs^d_\g(R_d)$;
\item $B(x_i,x_j)=-S\,(u_iv_j+u_jv_i)$;
\item the homomorphism $\varphi:k[\Cs^d_\g]\to R_d$ defined by $(x_1,\dots,x_d)$ satisfies $\varphi(\det T)=(-1)^d(d+1)!\,S^d\otimes u_1v_1\cdots u_dv_d$.
\end{enumerate}
\end{proposition}
\begin{proof}
(i) Moving the odd $u_i$ or $v_i$ past one odd generator,
\[
[x_i,x_j]=-\sum_{a,b}[e_a,e_b]\otimes\bigl(X^aX^b\,u_iu_j+X^aY^b\,u_iv_j+Y^aX^b\,v_iu_j+Y^aY^b\,v_iv_j\bigr).
\]
The four coefficients are $\{X,X\}$, $\{X,Y\}$, $\{Y,X\}=\{X,Y\}$ and $\{Y,Y\}$. For instance, in $U(\g)\otimes A$ one has $XX=\sum_{a,b}e_ae_bX^aX^b=\tfrac12\sum_{a,b}[e_a,e_b]X^aX^b$ (antisymmetrize in $a,b$ using $X^bX^a=-X^aX^b$), so $XX+XX=\sum_{a,b}[e_a,e_b]X^aX^b$; similarly $\sum_{a,b}[e_a,e_b]Y^aX^b=\sum_{a,b}[e_a,e_b]X^aY^b$. All four vanish in $A$.

(ii) With $B_{ab}=B(e_a,e_b)$,
\[
B(x_i,x_j)=-\sum_{a,b}B_{ab}\bigl(X^aX^b\,u_iu_j+X^aY^b\,u_iv_j+Y^aX^b\,v_iu_j+Y^aY^b\,v_iv_j\bigr).
\]
Symmetry of $B_{ab}$ and antisymmetry of $X^aX^b$, $Y^aY^b$ remove the first and last terms; $\sum B_{ab}Y^aX^b=-\sum B_{ab}X^bY^a=-S$ and $v_iu_j=-u_jv_i$ give $B(x_i,x_j)=-S\,u_iv_j-S\,u_jv_i$.

(iii) By (i) the evaluation $k[\g^d]\to R_d$ kills $I_d$. By (ii) and Lemma~\ref{lem:grass}, with $S$ even,
\[
\varphi(\det T)=\det\bigl(-S\,N_{ij}\bigr)=(-S)^d\det N=(-1)^d(d+1)!\,S^d\otimes u_1v_1\cdots u_dv_d .
\]
\end{proof}

\begin{corollary}\label{cor:nonzero}
If $d<\hv$, then $\det T\notin I_d$.
\end{corollary}
\begin{proof}
By External input~\ref{in:CDSW}(iii), $S^d\neq0$ in $A$ for $d\le\hv-1$. The algebra $A\otimes\Lambda$ is free over $A$ on the monomials in $u,v$, so $S^d\otimes u_1v_1\cdots u_dv_d\neq0$ and Proposition~\ref{prop:transfer}(iii) gives $\varphi(\det T)\neq0$. Since $\varphi$ factors through $k[\Cs^d_\g]=k[\g^d]/I_d$, the class of $\det T$ is nonzero.
\end{proof}

\section{The Gram component}\label{sec:gram}

\begin{lemma}[One $\Sg_n$-module for all $E$]\label{lem:polyfun}
For each $n\ge0$ there is a $G\times\Sg_n$-module $V_n$, independent of $E$, with natural isomorphisms
\[
\Oo(\Cs_E)_n\cong\bigl(E^{\otimes n}\otimes V_n\bigr)_{\Sg_n},
\]
where $\Sg_n$ permutes the tensor factors of $E^{\otimes n}$ with Koszul signs and those of $V_n$.
\end{lemma}
\begin{proof}
The supersymmetric power is $\Sym^n(E\otimes\g^*)=\bigl(E^{\otimes n}\otimes(\g^*)^{\otimes n}\bigr)_{\Sg_n}$. Let $c:\g^*\to(\g^*)^{\otimes2}$ be the transpose of the bracket and $K_2=c(\g^*)$, a space of antisymmetric tensors. The function $\rho_E(e\wedge e'\otimes\xi)$ is the class of $e\otimes e'\otimes c(\xi)$, so the degree-$n$ part of $J_E$ is the image of $E^{\otimes n}\otimes\bigl(K_2\otimes(\g^*)^{\otimes(n-2)}\bigr)$, which equals $\bigl(E^{\otimes n}\otimes K_n\bigr)_{\Sg_n}$ for the $\Sg_n$-submodule $K_n\subset(\g^*)^{\otimes n}$ generated by $K_2\otimes(\g^*)^{\otimes(n-2)}$. The functor of $\Sg_n$-coinvariants is exact in characteristic zero, hence $\Oo(\Cs_E)_n=(E^{\otimes n}\otimes V_n)_{\Sg_n}$ with $V_n=(\g^*)^{\otimes n}/K_n$.
\end{proof}

\begin{inputthm}[{Super Schur--Weyl duality, Sergeev \cite{Se85}, Berele--Regev \cite{BR87}}]\label{in:schur}
For a partition $\lambda\vdash n$ with Specht module $\mathrm{Sp}_\lambda$, the super Schur functor is $\Sf_\lambda(E)=(E^{\otimes n}\otimes\mathrm{Sp}_\lambda)_{\Sg_n}$. For $E=k^{m|p}$, $\Sf_\lambda(E)\neq0$ if and only if $\lambda_{m+1}\le p$, and the nonzero $\Sf_\lambda(E)$ are irreducible and pairwise non-isomorphic $\GL(E)$-modules. In particular $\Sf_\lambda(k^d)$ is the irreducible $\GL_d$-module of highest weight $\lambda$ (zero if $\lambda$ has more than $d$ parts), $\Sf_\lambda(\Pi k^e)\cong\Sf_{\lambda'}(k^e)$ as $\GL_e$-modules with $\lambda'$ the transpose partition, and $\Sym(E\otimes F)=\bigoplus_\lambda\Sf_\lambda(E)\otimes\Sf_\lambda(F)$ for $F$ purely even.
\end{inputthm}

Decompose $V_n=\bigoplus_{\lambda\vdash n}\mathrm{Sp}_\lambda\otimes M_\lambda$ with $M_\lambda=\Hom_{\Sg_n}(\mathrm{Sp}_\lambda,V_n)$, a $G$-module. Lemma~\ref{lem:polyfun} and External input~\ref{in:schur} give the $\GL(E)\times G$-decomposition
\begin{equation}\label{eq:decomp}
\Oo(\Cs_E)_n=\bigoplus_{\lambda\vdash n}\Sf_\lambda(E)\otimes M_\lambda,\qquad \Oo(\Cs_E)^G_n=\bigoplus_{\lambda\vdash n}\Sf_\lambda(E)\otimes M^G_\lambda .
\end{equation}

\begin{lemma}[Naturality of the Cartan restriction]\label{lem:natural}
For each $\lambda\vdash n$ there is a linear map $\phi_\lambda:M^G_\lambda\to\Sf_\lambda(\h^*)^W$, independent of $E$, with
\[
\Phi_E=\bigoplus_{\lambda}\mathrm{id}_{\Sf_\lambda(E)}\otimes\phi_\lambda .
\]
In particular $\Phi_E$ vanishes on $\Sf_\lambda(E)\otimes M^G_\lambda$ when $\lambda$ has more than $\rk\g$ parts.
\end{lemma}
\begin{proof}
Restriction $\g^*\to\h^*$ induces a $\Sg_n$-equivariant map $(\g^*)^{\otimes n}\to(\h^*)^{\otimes n}$, which kills $K_n$ because $\h$ is abelian; let $V_n\to(\h^*)^{\otimes n}$ be the induced map. Applying $(E^{\otimes n}\otimes-)_{\Sg_n}$ gives the restriction $\Oo(\Cs_E)_n\to\Sym^n(E\otimes\h^*)$, and $\Hom_{\Sg_n}(\mathrm{Sp}_\lambda,(\h^*)^{\otimes n})=\Sf_\lambda(\h^*)$ by External input~\ref{in:schur}. Taking $G$-invariants on the source, whose image is $W$-invariant, gives $\phi_\lambda$. Finally $\Sf_\lambda(\h^*)=0$ when $\lambda$ has more than $\dim\h$ parts.
\end{proof}

\begin{proof}[Proof of Theorem~\ref{thm:A}]
By \eqref{eq:decomp} and External input~\ref{in:schur}, the $\Sf_{(2^d)}(E)$-isotypic component of $\Oo(\Cs_E)^G_{2d}$ is $\Sf_{(2^d)}(E)\otimes L_d$ with $L_d=M^G_{(2^d)}$, which is independent of $E$.

\emph{Odd side.} For $E=\Pi k^2$, $\Sf_\lambda(\Pi k^2)=\Sf_{\lambda'}(k^2)$, and $\Sf_\mu(k^2)$ contains $\SL_2$-invariants exactly for $\mu=(c,c)$. In degree $2d$ this selects $\lambda'=(d,d)$, i.e.\ $\lambda=(2^d)$, and $\Sf_{(d,d)}(k^2)=\det^d$ is one-dimensional. Hence
\[
\bigl(A^G_{2d}\bigr)^{\SL_2}\cong L_d .
\]
By External input~\ref{in:CDSW}(i), $A^G=A^\g$ is generated by $S$ for connected $G$, and $S$ has degree $2$. Under $(X,Y)\mapsto(g_{11}X+g_{12}Y,\ g_{21}X+g_{22}Y)$ with $g\in\GL_2$ one has $S\mapsto\det(g)\,S$, because $B(X,X)=B(Y,Y)=0$ and $B(Y,X)=-B(X,Y)$; so $S$ is $\SL_2$-invariant. Therefore $(A^G_{2d})^{\SL_2}=k\,S^d$, and by (ii), (iii) its dimension is $1$ for $d<\hv$ and $0$ for $d\ge\hv$.

\emph{Even side.} For $E=k^d$, $\Sf_\lambda(k^d)$ with $\lambda\vdash2d$ contains $\SL_d$-invariants exactly for $\lambda=(2^d)$, where $\Sf_{(2^d)}(k^d)=\det^2$. So the $\SL_d$-invariant subspace of $k[\Cs^d_\g]^G_{2d}$ is $L_d$. It contains the class of $\det T$, which is nonzero for $d<\hv$ by Corollary~\ref{cor:nonzero} and therefore spans it.
\end{proof}

\section{The Cartan restriction and nilpotency}\label{sec:cartan}

\begin{lemma}[{Semisimple parts; cf.\ \cite{Ri88}, \cite[\S6]{Hu97}, \cite[\S3]{CN20}}]\label{lem:ss}
Let $\g$ be semisimple and $(x_1,\dots,x_d)\in\Cs^d_\g(k)$ with Jordan decompositions $x_i=s_i+n_i$. Then $s_1,\dots,s_d$ lie in one Cartan subalgebra and $f(x_1,\dots,x_d)=f(s_1,\dots,s_d)$ for every $f\in k[\g^d]^G$.
\end{lemma}
\begin{proof}
Each $\operatorname{ad}s_i$ and $\operatorname{ad}n_i$ is a polynomial in $\operatorname{ad}x_i$, so all $s_i,n_j$ commute pairwise. The $s_i$ span a toral subalgebra, contained in a Cartan subalgebra $\h'$. Let $L=Z_G(s_1,\dots,s_d)^\circ$, a connected reductive subgroup that fixes each $s_i$, with Lie algebra $\mathfrak l=\mathfrak z_\g(s_1,\dots,s_d)$. The commuting nilpotent elements $n_i\in\mathfrak l$ span an abelian subalgebra of nilpotent elements, which lies in the nilradical of a Borel subalgebra $\mathfrak b\subset\mathfrak l$. Let $B_L\supset T_L$ be a Borel subgroup and a maximal torus of $L$ with $\operatorname{Lie}B_L=\mathfrak b$, and take a cocharacter $\mu$ of $T_L$ with positive pairing with every root of the nilradical of $\mathfrak b$. Then $\mu(t)\cdot x_i=s_i+\mu(t)\cdot n_i\to s_i$ as $t\to0$, and $f(x)=f(\mu(t)\cdot x)\to f(s)$.
\end{proof}

\begin{proof}[Proof of Corollary~\ref{cor:B}]
By Corollary~\ref{cor:nonzero}, the class of $\det T$ in $k[\Cs^d_\g]^G$ is nonzero. Since $d>\dim\h$, vectors $x_1,\dots,x_d\in\h$ are linearly dependent, so $\Phi(\det T)=0$ (equivalently, Lemma~\ref{lem:natural} with $\lambda=(2^d)$). At every $k$-point of $\Cs^d_\g$, Lemma~\ref{lem:ss} gives $\det T(x)=\det T(s)$ with $s_1,\dots,s_d$ in a Cartan subalgebra of dimension $\rk\g<d$, so $\det T(x)=0$. By the Nullstellensatz, $\det T\in\sqrt{I_d}$, and its class is a nonzero nilpotent element of $k[\Cs^d_\g]^G$. Conjecture~3.1 of \cite{CN20} asserts that $\Phi$ is an isomorphism, and Conjecture~3.2 implies that $\Phi$ is injective \cite[\S3]{CN20}.
\end{proof}

\begin{remark}\label{rem:visible}
For $d\le\rk\g$ (and then $d<\hv$), $\Phi$ maps the one-dimensional space $L_d$ isomorphically onto $\Sf_{(2^d)}(\h^*)^W$, which is one-dimensional. Indeed, the Gram determinant on $\h^d$ is a nonzero element of $\Sf_{(2^d)}(\h^*)^W$. Conversely, by Steinberg's theorem \cite[Ch.~V, \S2, Exerc.~3]{Bou} the $W$-module $\Lambda^d\h^*$ is irreducible for $0\le d\le\rk\g$; it is defined over $\mathbb R$, hence self-dual and orthogonal, so $\dim\Sym^2(\Lambda^d\h^*)^W=1$. Since $\Sf_{(2^d)}(\h^*)$ is a $\GL(\h^*)$-summand of $\Sym^2(\Lambda^d\h^*)$, $\dim\Sf_{(2^d)}(\h^*)^W\le1$.
\end{remark}

\section{Two examples}\label{sec:examples}

For a multilinear function $m$ of six matrix arguments $(a_1,a_2,a_3,b_1,b_2,b_3)$ and $x_1,x_2,x_3\in\g$ write
\[
\langle m\rangle=\sum_{\sigma,\tau\in\Sg_3}\sgn\sigma\,\sgn\tau\;m\bigl(x_{\sigma(1)},x_{\sigma(2)},x_{\sigma(3)},x_{\tau(1)},x_{\tau(2)},x_{\tau(3)}\bigr).
\]
For a representation of $\g$ with trace $\operatorname{tr}$ and $B$ its trace form, $\langle\operatorname{tr}(a_1b_1)\operatorname{tr}(a_2b_2)\operatorname{tr}(a_3b_3)\rangle=6\det T$.

\begin{exampleblock}[$\g_2$, $d=3$]\label{ex:g2}
Here $\rk\g_2=2<3<4=\hv$. Let $Z_3=(\Sf_{(2,2,2)}\g_2)^{G_2}$ be the $\SL_3$-invariant part of $k[\g_2^3]^{G_2}$ of degree $6$, and $Y_3=(\g_2\otimes\Sf_{(2,1,1)}\g_2)^{G_2}$. Since $G_2\times\SL_3$ is reductive, $I_3\cap Z_3$ is the image of the $G_2\times\SL_3$-invariants of $\g_2^*\otimes\Lambda^2k^3\otimes\Sym^4(k^3\otimes\g_2^*)$ under multiplication by the commutator coordinates, and Pieri's rule ($\Lambda^2k^3\otimes\Sf_\mu k^3$ contains $\det^2$ iff $\mu=(2,1,1)$) identifies that space with $Y_3$. Exact character computations give
\[
\dim Z_3=4,\qquad \dim Y_3=3 ,
\]
so $Z_3/(I_3\cap Z_3)\neq0$ by counting alone, Lemma~\ref{lem:natural} gives $\Phi(Z_3)=0$, and Theorem~\ref{thm:A} identifies the quotient with $k\cdot\det T$. With traces in the $7$-dimensional representation, let
\begin{gather*}
S_{6a}=\langle\operatorname{tr}(a_1a_2a_3b_1b_2b_3)\rangle,\qquad S_{6b}=\langle\operatorname{tr}(a_1b_1a_2b_2a_3b_3)\rangle,\\
F_{42}=\langle\operatorname{tr}(a_1a_2b_1b_2)\operatorname{tr}(a_3b_3)\rangle,\qquad F_{33}=\langle\operatorname{tr}(a_1a_2a_3)\operatorname{tr}(b_1b_2b_3)\rangle .
\end{gather*}
Then $Z_3$ has basis $S_{6a},S_{6b},F_{42},\det T$; $I_3\cap Z_3$ has basis $S_{6a},S_{6b},F_{42}$; and $-18S_{6a}-6S_{6b}+3F_{42}+4F_{33}=0$ (Appendix~\ref{app:comp}). For $\g_2$ and $d=3$, $\Phi$ is injective on every $\GL_3$-isotypic component of degree at most $10$ except $\Sf_{(2,2,2)}$, where its kernel is $k\cdot\det T$, and $(\det T)^2\in I_3$ (Appendix~\ref{app:comp}).

For comparison, $\mathfrak{sl}_3$, $\mathfrak{so}_5$ and $\mathfrak{sp}_4$ have $\rk=2$ and $\hv=3$, so Theorem~\ref{thm:A} gives $\det T\in I_3$. Explicitly, with traces in the defining representation,
\begin{align*}
\mathfrak{sl}_3:\ \det T&=\Bigl\langle-\tfrac56\operatorname{tr}([a_1,a_2]a_3b_1b_2b_3)-\tfrac23\operatorname{tr}([a_1,a_2]b_1a_3b_2b_3)\\
&\qquad+\tfrac5{18}\operatorname{tr}([a_1,a_2]a_3)\operatorname{tr}(b_1b_2b_3)\Bigr\rangle,\\
\mathfrak{sp}_4:\ \det T&=\Bigl\langle-\tfrac16\operatorname{tr}([a_1,a_2]a_3b_1b_2b_3)-\tfrac56\operatorname{tr}([a_1,a_2]b_1a_3b_2b_3)\\
&\qquad+\tfrac14\operatorname{tr}([a_1,a_2]b_1b_3)\operatorname{tr}(a_3b_2)\Bigr\rangle,\\
\mathfrak{so}_5:\ \det T&=\Bigl\langle\tfrac{10}3\operatorname{tr}([a_1,a_2]a_3b_1b_2b_3)-\tfrac43\operatorname{tr}([a_1,a_2]b_1a_3b_2b_3)\\
&\qquad-\tfrac23\operatorname{tr}([a_1,a_2]a_3)\operatorname{tr}(b_1b_2b_3)+\operatorname{tr}([a_1,a_2]b_1b_3)\operatorname{tr}(a_3b_2)\Bigr\rangle .
\end{align*}
Every term contains a commutator, so each right-hand side lies in $I_3$.
\end{exampleblock}

\begin{exampleblock}[$\mathfrak{so}_7$, $d=4$]\label{ex:so7}
Here $\rk=3<4<5=\hv$, and Corollary~\ref{cor:B} gives a nonzero nilpotent class $\det T\in k[\Cs^4_{\mathfrak{so}_7}]^{SO_7}$ with $\Phi(\det T)=0$. The determinant $\det T$ is $O_7$-invariant, so the same element lies in $k[\Cs^4_{\mathfrak{so}_7}]^{O_7}$. In the $\SL_4$-invariant part of degree $8$,
\[
\dim\bigl(\Sf_{(2,2,2,2)}\mathfrak{so}_7\bigr)^{SO_7}=9,\qquad \dim\bigl(\mathfrak{so}_7\otimes\Sf_{(2,2,1,1)}\mathfrak{so}_7\bigr)^{SO_7}=12,
\]
and the commutator substitution has rank at least $8$ on the second space (rank $8$ at random points modulo three primes; Appendix~\ref{app:comp}). Theorem~\ref{thm:A} gives rank exactly $8$.
\end{exampleblock}

\section{The orthogonal case of Song--Xia--Xu}\label{sec:sxx}

For $G=O_n$ and for $G=SO_n$ with $n$ odd, \cite[Thm.~1.2]{SXX23} is \cite[Thm.~3.4]{SXX23}, deduced from \cite[Prop.~3.1, Lemmas~3.2, 3.3]{SXX23}; for the orthogonal groups the degree bound \cite[Lemma~3.3]{SXX23} is deduced in \cite[\S5.2]{SXX23} from \cite[Lemma~5.3]{SXX23}, whose proof invokes \cite[Lemma~5.2]{SXX23} for the same $d$. For $SO_{2m}$ the proof of \cite[Thm.~4.3]{SXX23} uses \cite[Thm.~3.4]{SXX23} for $O_{2m}$. The symplectic case rests on the Pfaffian of Chen--Ng\^o and is independent of \cite[Lemma~5.2]{SXX23}. In this section $(t_{ij})$ is the generic skew-symmetric $(2d+3)\times(2d+3)$ matrix (called $T$ in \cite{SXX23}), $h_{i_1\cdots i_k}$ is the Pfaffian of the matrix obtained by deleting the rows and columns $i_1,\dots,i_k$ (a repeated index leaves a matrix of odd size, so $h_{i,i,\dots}=0$, consistent with $\Asxx(1,1)=2h^2_{1,3,4}$ in the case $d=1$ of \cite[\S5.2]{SXX23}), and $\Asxx$ is the $2d\times2d$ matrix of \cite[\S5.2]{SXX23} (called $A$ there),
\[
\Asxx(i,j)=h^2_{i,j,2d+2}+h^2_{i,2d+1,2d+2}+h^2_{j,2d+1,2d+2},\qquad 1\le i,j\le2d .
\]
\cite[Lemma~5.2]{SXX23} asserts that $h_{2d+2}=0$ and $h_{2d+3}\det\Asxx\neq0$ have a common solution.

\begin{proposition}\label{prop:rank3}
On the hypersurface $h_{2d+2}=0$, the matrix $\Asxx$ has rank at most $3$. Consequently $\det\Asxx\in(h_{2d+2})$ for every $d\ge2$, and \cite[Lemma~5.2]{SXX23} fails for every $d\ge2$.
\end{proposition}
\begin{proof}
Let $P$ be the skew-symmetric matrix of size $2d+2$ obtained from $(t_{ij})$ by deleting row and column $2d+2$, so $\Pf P=h_{2d+2}$; its rows and columns $1,\dots,2d+1$ are those of $(t_{ij})$, and its last one is the original $2d+3$. For $1\le a<b\le2d+2$ put $P^\sharp_{ab}=(-1)^{a+b}\Pf P(\hat a,\hat b)$ and $P^\sharp_{ba}=-P^\sharp_{ab}$. Expansion of the Pfaffian along a row gives $PP^\sharp=(\Pf P)\cdot\mathbf 1$, and $h_{i,j,2d+2}=\pm P^\sharp_{ij}$, $h_{i,2d+1,2d+2}=\pm P^\sharp_{i,2d+1}$ for $i,j\le2d$.

On $\Pf P=0$ the columns of $P^\sharp$ lie in $\ker P$. If $\rk P=2d$ this kernel has dimension $2$; if $\rk P\le2d-2$, every principal submatrix of size $2d$ is singular and $P^\sharp=0$. In both cases the skew matrix $P^\sharp$ has rank at most $2$, hence $P^\sharp=vw^{\mathsf T}-wv^{\mathsf T}$ for vectors $v,w$. Squares remove the signs:
\[
h^2_{i,j,2d+2}=(v_iw_j-v_jw_i)^2,\qquad h^2_{i,2d+1,2d+2}=\ell_i^2,\quad \ell_i=w_{2d+1}v_i-v_{2d+1}w_i .
\]
With $f_i=(v_i^2,\,w_i^2,\,v_iw_i,\,1)$ and $a=w_{2d+1}^2$, $b=v_{2d+1}^2$, $c=-2w_{2d+1}v_{2d+1}$, expanding the squares gives $\Asxx(i,j)=f_iMf_j^{\mathsf T}$ with
\[
M=\begin{pmatrix}0&1&0&a\\1&0&0&b\\0&0&-2&c\\a&b&c&0\end{pmatrix},\qquad
\det M=-2\det\begin{pmatrix}0&1&a\\1&0&b\\a&b&0\end{pmatrix}+c^2=-4ab+c^2=0 .
\]
Thus $\rk\Asxx\le\rk M\le3<2d$, and $\det\Asxx$ vanishes at every $k$-point of $h_{2d+2}=0$. The Pfaffian $h_{2d+2}$ is irreducible \cite[Lemma~B.2.10]{GW09}, so $(h_{2d+2})$ is prime and the Nullstellensatz gives $\det\Asxx\in(h_{2d+2})$.
\end{proof}

\begin{remark}\label{rem:charp}
The proof uses polynomial identities over $\mathbb Z$ and linear algebra over a field of characteristic $\neq2$. Hence \cite[Lemma~5.2]{SXX23} also fails in characteristic $p>2$ for $d\ge2$, and the proofs of \cite[Thms.~3.5, 4.4]{SXX23} (Theorem~1.2(2) there) for the orthogonal groups are incomplete for $d\ge2$.
\end{remark}

\begin{remark}[The failing step]\label{rem:step}
The induction in the proof of \cite[Lemma~5.2]{SXX23} (arXiv v3, p.~17) asserts that $\det\Asxx$ has degree at most $4d-4$ in $t_{12}$, with coefficient of $t_{12}^{4d-4}$ equal to
\[
\det\begin{pmatrix}2h^2_{1,2d+1,2d+2}&h^2_{1,2,2d+2}+h^2_{1,2d+1,2d+2}+h^2_{2,2d+1,2d+2}\\ h^2_{1,2,2d+2}+h^2_{1,2d+1,2d+2}+h^2_{2,2d+1,2d+2}&2h^2_{2,2d+1,2d+2}\end{pmatrix}\cdot\det A' ,
\]
with $A'$ the matrix of the same form for the matrix obtained by deleting rows and columns $1,2$. For $j\ge3$ the entries $\Asxx(1,j)$ and $\Asxx(2,j)$ contain $h^2_{j,2d+1,2d+2}$, which has degree $2$ in $t_{12}$. The leading $t_{12}^2$-parts of rows $1$ and $2$ coincide, so the coefficient of $t_{12}^{4d}$ vanishes; at $d=2$ and $d=3$ the exact degree is $4d-2$ ($6$ and $10$), and at random integer values of the other $t_{ij}$ the coefficient of $t_{12}^{4d-4}$ differs from the displayed product (Appendix~\ref{app:comp}). The base case $d=1$ is correct: there $\det\Asxx=-(t_{35}^2+(t_{25}+t_{15})^2)(t_{35}^2+(t_{25}-t_{15})^2)$ (arXiv v3, p.~16), and Proposition~\ref{prop:rank3} holds trivially because $2d=2\le3$.
\end{remark}

\begin{corollary}\label{cor:SXX}
Let $n=2m+1\ge7$ with $m+1\le d\le2m-2$, or $n=2m\ge8$ with $m+1\le d\le2m-3$. For $G=O_n$ and $G=SO_n$, with $W=N_G(\h)/Z_G(\h)$, the restriction $\Phi:k[\Cs^d_{\mathfrak{so}_n}]^G\to k[\h^d]^W$ has a nonzero kernel and $k[\Cs^d_{\mathfrak{so}_n}]^G$ has a nonzero nilradical. In particular \cite[Thm.~1.2(1), Cor.~1.3]{SXX23} fail for $n=7$, $d=4$. Since \cite[Prop.~3.1, Lemma~3.2]{SXX23} are independent of \cite[\S5.2]{SXX23} and \cite[Thm.~3.4]{SXX23} follows formally from them and \cite[Lemma~3.3]{SXX23}, the degree bound \cite[Lemma~3.3]{SXX23} and \cite[Lemma~5.3]{SXX23}, which implies it, are false for $G=O_n$ and for $G=SO_n$ with $n$ odd in the same range.
\end{corollary}
\begin{proof}
These are the pairs $(B_m,d)$ and $(D_m,d)$ of Table~\ref{tab:windows}. Corollary~\ref{cor:B} gives the statement for $SO_n$; $\det T$ is $O_n$-invariant, so the same element works for $O_n$.
\end{proof}

\section{Super-commuting schemes and $\mathcal N=4$ Yang--Mills theory}\label{sec:N4}

Theorem~\ref{thm:A} holds for every finite-dimensional super vector space $E$. For $E=k^{3|2}$ the scheme $\Cs_E=\Cs^{3|2}_\g$ is the super-commuting scheme of \cite[\S2]{Ch26}. By Chang--Yin \cite{CY13}, the classical $Q$-cohomology of $\mathcal N=4$ super-Yang--Mills theory with gauge algebra $\g$ is the relative Lie algebra cohomology $H^\bullet(\g[\mathsf A],\g;\mathbb C)$ with $\mathsf A=\mathbb C[z_+,z_-]\otimes\Lambda(\theta_1,\theta_2,\theta_3)$. By \cite[Prop.~4]{Ch26}, its part of top degree $p=|n|$ and multidegree $n=(n_{z_+},n_{z_-},n_{\theta_1},n_{\theta_2},n_{\theta_3})$ is $\Oo(\Cs^{3|2}_\g)^G_n$, the three even coordinates being $\partial_{\theta_i}\Psi|_{z=\theta=0}$ and the two odd ones $\partial_{z_\pm}\Psi|_{z=\theta=0}$ for the BPS superfield $\Psi$ of \cite[\S3]{Ch26}. Following \cite[Def.~1]{Ch26}, a non-Cartan class is a nonzero element of $\ker\Phi_{k^{3|2}}$; the level of $n$ is $3n_{z_+}+3n_{z_-}+2(n_{\theta_1}+n_{\theta_2}+n_{\theta_3})$.

\begin{corollary}\label{cor:N4}
Let $\g$ be simple and $1\le d<\hv$. With respect to the Borel subalgebra of $\mathfrak{gl}(3|2)$ with the even basis vectors first, $\Sf_{(2^d)}(k^{3|2})$ is nonzero with highest weight
\[
\bigl((2^{\min(d,3)},0^{3-\min(d,3)})\,\big|\,(c,c)\bigr),\qquad c=\max(d-3,0).
\] Hence $\Oo(\Cs^{3|2}_\g)^G$ contains $\Sf_{(2^d)}(k^{3|2})\otimes L_d\neq0$, of degree $2d$; its highest-weight vector has multidegree
\[
n=(0,0,2,0,0),\ (0,0,2,2,0),\ (d-3,d-3,2,2,2)\qquad (d=1,\ d=2,\ d\ge3),\qquad p=|n|=2d,
\]
of level $4d+2c$. Its part of multidegree $n=(d,d,0,0,0)$ is $k\,S^d$, with $S=B(X,Y)$ for the two odd coordinates $X=\partial_{z_+}\Psi|_{z=\theta=0}$, $Y=\partial_{z_-}\Psi|_{z=\theta=0}$. For $d>\rk\g$ this component lies in $\ker\Phi_{k^{3|2}}$.
\end{corollary}
\begin{proof}
The hook condition $\lambda_4\le2$ of External input~\ref{in:schur} holds for $\lambda=(2^d)$. For the distinguished Borel subalgebra the highest weight of $\Sf_\lambda(k^{3|2})$ is $\bigl(\lambda_1,\lambda_2,\lambda_3\,\big|\,\max(\lambda'_1-3,0),\max(\lambda'_2-3,0)\bigr)$ \cite{BR87}, which for $\lambda=(2^d)$ is the stated weight. Write $k^{3|2}=k^3\oplus\Pi k^2$. In Lemma~\ref{lem:polyfun} the part of $E^{\otimes n}$ of degree zero in $k^3$ is $(\Pi k^2)^{\otimes n}$, so the part of $\Oo(\Cs_{k^{3|2}})$ of degree zero in the even coordinates is $\Oo(\Cs_{\Pi k^2})=A$, and the corresponding part of $\Sf_{(2^d)}(k^{3|2})$ is $\Sf_{(2^d)}(\Pi k^2)=\det^d$, of multidegree $(d,d,0,0,0)$. By the odd side of the proof of Theorem~\ref{thm:A}, the part of the component of multidegree $(d,d,0,0,0)$ is $(A^G_{2d})^{\SL_2}=k\,S^d$. The rest follows from Theorem~\ref{thm:A} and Lemma~\ref{lem:natural}.
\end{proof}

\begin{corollary}\label{cor:spE}
If $\Phi_{k^d}$ is an isomorphism for every $d\ge1$, then $\Phi_E$ is an isomorphism for every finite-dimensional super vector space $E$. In particular, for $\g=\mathfrak{sp}_{2N}$ the map $\Phi_{k^{3|2}}$ is an isomorphism, and every nonzero element of $\Oo(\Cs^{3|2}_{\mathfrak{sp}_{2N}})^G$ has nonzero Cartan restriction.
\end{corollary}
\begin{proof}
By Lemma~\ref{lem:natural}, $\Phi_{k^d}=\bigoplus_\lambda\mathrm{id}\otimes\phi_\lambda$ with $\Sf_\lambda(k^d)\neq0$ for every $\lambda$ with at most $d$ parts; so all $\phi_\lambda$ are isomorphisms, and $\Phi_E=\bigoplus_\lambda\mathrm{id}\otimes\phi_\lambda$ is an isomorphism. For $\mathfrak{sp}_{2N}$ the hypothesis is \cite[Thm.~4.1]{CN21}.
\end{proof}

\begin{remark}\label{rem:N4}
\begin{enumerate}[label=(\alph*),leftmargin=*]
\item For $\g_2$ ($\rk=2$, $\hv=4$), $d=3$ gives the non-Cartan class of $\det T$ at $n=(0,0,2,2,2)$, of level $12$; in the conventions of \cite[\S5.2]{Ch26} its charges are $(J_1,J_2,q_1,q_2,q_3)=(0,0,2,2,2)$. It is a purely scalar element of the classical $Q$-cohomology.
\item For $\mathfrak{so}_7$ ($\rk=3$, $\hv=5$), $d=4$ gives a non-Cartan class at $n=(1,1,2,2,2)$, $p=8$, of level $18$. By \cite[\S3, Table~1 and Discussions]{CL25}, $\Oo(\Cs^{3|2}_{\mathfrak{so}_7})^{SO_7}_n$ is $85$-dimensional there ($84$ for $\mathfrak{sp}_6$), and super-abelianization, which is $\Phi_{k^{3|2}}$, has a one-dimensional kernel on it. The class of Corollary~\ref{cor:N4} therefore spans this kernel, and the non-Cartan class of \cite[Thm.~1]{Ch26} is a nonzero multiple of it. A representative of that class, due to Choi and Lee \cite{CE25} and written in superfield form in \cite[eq.~(16)]{Ch26}, is $\varepsilon_{ijk}\varepsilon_{lmn}Z_{il}\operatorname{Tr}(\Psi_j\Psi_m)\operatorname{Tr}(\Psi_k\Psi_n)$ with $Z_{il}=\operatorname{Tr}(\Psi_i\Psi_l)\operatorname{Tr}(\Psi_{\dot\alpha}\Psi^{\dot\alpha})-\operatorname{Tr}(\Psi_i\Psi_{\dot\alpha})\operatorname{Tr}(\Psi_l\Psi^{\dot\alpha})-4\operatorname{Tr}(\Psi_{\dot\alpha}\Psi^{\dot\alpha}\Psi_{(i}\Psi_{l)})$; in the notation of \cite{Ch26}, $\Psi_i$ and $\Psi_{\dot\alpha}$ are the even and odd coordinates.
\item The Langlands-dual pair $(\mathfrak{so}_{2N+1},\mathfrak{sp}_{2N})$ has rank $N$ and dual Coxeter numbers $2N-1$ and $N+1$. Corollary~\ref{cor:N4} gives non-Cartan classes for $\mathfrak{so}_{2N+1}$ at $n=(d-3,d-3,2,2,2)$ for $N+1\le d\le2N-2$, and by Corollary~\ref{cor:spE} the Cartan restriction is injective for $\mathfrak{sp}_{2N}$. For $N=3$ this is the non-Cartan part of the classical mismatch at level $18$ in \cite{CL25}; for $N=4$ it gives non-Cartan $\mathfrak{so}_9$ classes at $n=(2,2,2,2,2)$ and $(3,3,2,2,2)$.
\item In the notation of \cite[\S2]{KvLM26} the odd coordinates $X,Y$ are the gaugino components $\lambda_\pm$, and $S$ is a nonzero multiple of $\operatorname{Tr}(\lambda_+\lambda_-)$. For $\rk\g<d<\hv$ the vector $S^d$ of Corollary~\ref{cor:N4} is the non-Coulomb monotone class $[S^d]$ of \cite[eq.~(1)]{KvLM26}. Corollary~\ref{cor:N4} places $S^d$ and the class at $n=(d-3,d-3,2,2,2)$ in one irreducible $\mathfrak{gl}(3|2)$-module. For $\mathfrak{so}_7$, $d=4$, this agrees with \cite[\S5, App.~D]{KvLM26}, where $S^4$ is shown to be a superconformal descendant of the class in~(b).
\end{enumerate}
\end{remark}

\appendix
\section{Computations}\label{app:comp}

Exact counts come from Weyl characters (two independent character routes for the $d=3$ rows). Ranks are ranks of matrices of function values at random points: such a rank is a lower bound for the rank over $\mathbb Q$, and it is exact when it equals the dimension of the source or the target. The $d=3$ ranks were also computed exactly over $\mathbb Z$, with Example~\ref{ex:g2} recomputed in a second realization of $\g_2$ (the stabilizer in $\mathfrak{so}_7$ of the octonionic $3$-form).

\begin{table}[ht]
\centering\small
\setlength{\tabcolsep}{4pt}
\begin{tabular}{lcccccc}
\toprule
$\g$, $d$ & $(\dim Y_d,\dim Z_d)$ & rank & $\det T\in I_d$ & $\dim Z_d-$rank & $\dim L_d$ & $\Phi(\det T)$\\
\midrule
$\mathfrak{sl}_3$, $3$ & $(4,3)$ & $3$ & yes & $0$ & $0$ & --\\
$\mathfrak{so}_5,\mathfrak{sp}_4$, $3$ & $(4,4)$ & $4$ & yes & $0$ & $0$ & --\\
$\g_2$, $3$ & $(3,4)$ & $3$ & no & $1$ & $1$ & $0$\\
$\mathfrak{so}_7,\mathfrak{sp}_6$, $3$ & $(4,5)$ & $4$ & no & $1$ & $1$ & $\neq0$\\
$\mathfrak{sl}_4$, $3$ & $(5,5)$ & $\ge4$ & no & $\le1$ & $1$ & $\neq0$\\
$\mathfrak{so}_5$, $4$ & $(9,5)$ & $5$ & yes & $0$ & $0$ & --\\
$\g_2$, $4$ & $(7,5)$ & $5$ & yes & $0$ & $0$ & --\\
$\mathfrak{sl}_4$, $4$ & $(20,8)$ & $8$ & yes & $0$ & $0$ & --\\
$\mathfrak{sp}_6$, $4$ & $(12,8)$ & $8$ & yes & $0$ & $0$ & --\\
$\mathfrak{so}_7$, $4$ & $(12,9)$ & $\ge8$ & no & $\le1$ & $1$ & $0$\\
\bottomrule
\end{tabular}
\caption{$Z_d=(\Sf_{(2^d)}\g)^G$ and $Y_d=(\g\otimes\Sf_{(2^{d-2},1,1)}\g)^G$, the source of the commutator substitution; ``rank'' is the rank of $Y_d\to Z_d$ and $\dim L_d$ is given by Theorem~\ref{thm:A}. Entries ``no'' in the column $\det T\in I_d$ follow from Corollary~\ref{cor:nonzero} for the rows where the computed rank is a lower bound ($\mathfrak{sl}_4$, $d=3$ and $\mathfrak{so}_7$, $d=4$) and from the exact computation in the other rows.}\label{tab:comp}
\end{table}

\begin{itemize}[leftmargin=*]
\item $\g_2$, $d=3$: the $265$ contractions $\langle m\rangle$ of fixed-point-free multilinear trace monomials (in the $7$-dimensional representation) have rank $4=\dim Z_3$ over $\mathbb Z$ at $16$ random integer points; adding $150$ elements of $Z_3$ built from the $3$-form $\varphi$ and the $4$-form ${*}\varphi$ leaves the rank at $4$. The contractions $\langle m\rangle$ of five-slot monomials with a commutator in one slot span a space of rank $3=\dim Y_3$, containing $S_{6a},S_{6b},F_{42},F_{33}$; adding $\det T$ raises the rank to $4$.
\item $\g_2$, degree $12$: the $\SL_3$-invariant part of $k[\g_2^3]^{G_2}$ of degree $12$ is $(\Sf_{(4,4,4)}\g_2)^{G_2}$, of dimension $30$; the commutator image has rank $30$ modulo $10^9+7$, so $(\det T)^2\in I_3$.
\item $\g_2$, $d=3$, degree at most $10$: for every partition $\lambda$ with $|\lambda|\le10$ and at most $3$ parts, $\Phi$ is injective on the $\Sf_\lambda$-isotypic component of $k[\Cs^3_{\g_2}]^{G_2}$, except for $\lambda=(2,2,2)$.
\item Lemma~\ref{lem:grass}: checked symbolically for $d\le5$.
\item Proposition~\ref{prop:rank3}: at random rational points of $h_{2d+2}=0$, $\rk\Asxx=2$ and $\det\Asxx\neq0$ for $d=1$; $\rk\Asxx=3$ and $\det\Asxx=0$ for $d=2,3,4$. Remark~\ref{rem:step}: at $d=2$ and $d=3$, $\deg_{t_{12}}\det\Asxx=4d-2$, and the coefficient of $t_{12}^{4d-4}$ differs from the product in Remark~\ref{rem:step} at three random integer points.
\end{itemize}

\begin{thebibliography}{99}\small
\bibitem{BR87} A.~Berele and A.~Regev, Hook Young diagrams with applications to combinatorics and to representations of Lie superalgebras, \emph{Adv.\ Math.} 64 (1987), 118--175.
\bibitem{Bou} N.~Bourbaki, \emph{Groupes et alg\`ebres de Lie}, Chapitres~IV--VI, Hermann, 1968.
\bibitem{CDSW} F.~Cachazo, M.~R.~Douglas, N.~Seiberg and E.~Witten, Chiral rings and anomalies in supersymmetric gauge theory, \emph{JHEP} 12 (2002) 071, arXiv:hep-th/0211170.
\bibitem{CF23} M.~Cederwall and G.~Ferretti, The chiral ring of $D=4$, $\mathcal N=1$ SYM with exceptional gauge groups, arXiv:2311.12119.
\bibitem{Ch26} C.-M.~Chang, Super-Chevalley restriction and relative Lie algebra cohomology over the $2|3$ algebra, arXiv:2604.23549.
\bibitem{CL25} C.-M.~Chang and Y.-H.~Lin, Violation of S-duality in classical $Q$-cohomology, arXiv:2510.24008.
\bibitem{CY13} C.-M.~Chang and X.~Yin, 1/16 BPS states in $\mathcal N=4$ super-Yang--Mills theory, arXiv:1305.6314.
\bibitem{CN20} T.-H.~Chen and B.~C.~Ng\^o, On the Hitchin morphism for higher-dimensional varieties, \emph{Duke Math.\ J.} 169 (2020), 1971--2004, arXiv:1905.04741.
\bibitem{CN21} T.-H.~Chen and B.~C.~Ng\^o, Invariant theory for the commuting scheme of symplectic Lie algebras, arXiv:2102.01849.
\bibitem{CE25} J.~Choi and E.~Lee, Konishi lifts a black hole, arXiv:2511.09519.
\bibitem{Do09} M.~Domokos, Vector invariants of a class of pseudoreflection groups and multisymmetric syzygies, \emph{J.\ Lie Theory} 19 (2009), 507--525.
\bibitem{EK03} P.~Etingof and V.~Kac, On the Cachazo--Douglas--Seiberg--Witten conjecture for simple Lie algebras, arXiv:math/0305175.
\bibitem{Et03} P.~Etingof, On the Cachazo--Douglas--Seiberg--Witten conjecture for simple Lie algebras, II, arXiv:math/0312452.
\bibitem{GW09} R.~Goodman and N.~R.~Wallach, \emph{Symmetry, Representations, and Invariants}, Graduate Texts in Mathematics 255, Springer, 2009.
\bibitem{Go26} P.~Goyal, The commuting scheme of $\g_2$, arXiv:2610.05018.
\bibitem{Hu97} M.~Hunziker, Classical invariant theory for finite reflection groups, \emph{Transform.\ Groups} 2 (1997), 147--163.
\bibitem{KvLM26} T.~Kibe, S.~van~Leuven and K.~Mathieson, Non-Coulomb $Q$-cohomology and quantum lifting, arXiv:2610.22332.
\bibitem{Ku08} S.~Kumar, On the Cachazo--Douglas--Seiberg--Witten conjecture for simple Lie algebras, \emph{J.\ Amer.\ Math.\ Soc.} 21 (2008), 797--808, arXiv:math/0412201.
\bibitem{LNY24} P.~Li, D.~Nadler and Z.~Yun, Functions on the commuting stack via Langlands duality, \emph{Ann.\ of Math.} 200 (2024), 609--748, arXiv:2301.02618.
\bibitem{Ri88} R.~W.~Richardson, Conjugacy classes of $n$-tuples in Lie algebras and algebraic groups, \emph{Duke Math.\ J.} 57 (1988), 1--35.
\bibitem{Se85} A.~N.~Sergeev, The tensor algebra of the identity representation as a module over the Lie superalgebras $\mathfrak{gl}(n,m)$ and $Q(n)$, \emph{Math.\ USSR Sb.} 51 (1985), 419--427.
\bibitem{SY26} J.~Sempliner and X.~Yang, On the normality of the commuting scheme, arXiv:2609.27360.
\bibitem{SXX23} L.~Song, X.~Xia and J.~Xu, A higher-dimensional Chevalley restriction theorem for orthogonal groups, \emph{Adv.\ Math.} 426 (2023), 109104, arXiv:2207.03147.
\bibitem{Va07} F.~Vaccarino, Linear representations, symmetric products and the commuting scheme, \emph{J.\ Algebra} 317 (2007), 634--641.
\bibitem{Wi03} E.~Witten, Chiral ring of $Sp(N)$ and $SO(N)$ supersymmetric gauge theory in four dimensions, arXiv:hep-th/0302194.
\end{thebibliography}

\end{document}
